\documentclass{article}
\usepackage{amsmath,hyperref}
\title{Voluntary preparation when a fallback option becomes locally exercisable}
\date{}
\begin{document}
\maketitle

\paragraph{Scope and timing.}
Two organizations operating civil inspection crews can arrange an optional joint service. The requesting crew A values the service at $v$ and has to prepare at irreversible cost $c$ before the local conditions are known. Crew B owns the other necessary resource, incurs $k$ if the service occurs, and can quote a fee in transferable service credits after preparation. Credit transfers are enforceable on actual performance and no advance price commitment is initially admissible. Both crews can decline trade. If trade does not happen, A can use a licensed, physically independent fallback worth $a>0$ only in some histories, and otherwise gets zero. It cannot obtain both benefits. Assume $V=v-k>a$ and nonnegative values and costs. The fallback is optional useful work, not a waiver of any public safety duty.

The source histories have a fixed declared distribution. A fallback grant is new authority, not proof that the apparatus is free. At offer cut 4 the local actor must have recognized the grant and its valid issuer, received evidence in time, and be able to begin a two-tick operation no earlier than actual dependency release, finishing by the grant expiry and by tick 6. A genuinely clear, timely early grant passes; a late copy, still-occupied source, expired or unrecognized grant does not. No inspection is launched from the latest remote record alone. For the economic comparison, the local decision-maker receives valid evidence of which option is actually exercisable; the price-setting role may or may not receive it before setting terms. The order matters even though every observation and all authorized disclosures are truthful. The prototype assumes these respective disclosures can be enforced by declared source-owned channels. It does not implement authenticated field networking. Service feasibility and payment verification are conditioned on a separate timely safety certificate in these payoff calculations.

\paragraph{Backwards induction at the fee cut.}
Let $q$ be the probability that the fallback is usable. With the two types indistinguishable to the fee setter, the only two potentially optimal offers are pooling at $p=v-a$ or screening at $p=v$. The former trades in both types and pays B $V-a$ while A earns $a$ gross in either type; the latter earns B $(1-q)V$ and A $qa$, with actual joint trade in the no-option type only. We select pooling at exact seller-profit equality; later we use strict margins where conclusions must not depend on that convention. Thus pooling occurs when $qV\ge a$. If the fee setter knows the actual option status, it can demand each type's maximum willingness to pay, trade in both, and leave A only $qa$ before its setup costs. The price mechanism uses acceptance at the reservation value; a strict perturbation feasible with fine credits preserves the strict ex ante examples, not every exact equality.

The threshold is not an artifact of testing only posted prices. For a general ex-post incentive-compatible direct lottery in this single-unit game, write $x_0,x_1$ for the probabilities of joint trade for types with zero and positive outside options, respectively. Utilities $u_0,u_1$ are in excess of the rejection allocation. Incentive compatibility gives $u_0\ge u_1+a x_1$, and participation gives $u_1\ge0$. Seller revenue net of physical performance cost is no larger than $(1-q)V+(qV-a)x_1$. The bound is attained by the two prices above, respectively $x_1=0,1$. At the profit tie, lotteries may provide other preparation incentives; the exact fixture below does not rely on this. This is standard single-seller screening under outside values, not a novelty claim about auction design. The derivation also allows state-contingent offers at a later cut only if that information is actually lawful to the price setter.

For example take $v=18,k=4,a=6,c=4$ and a seller-side fixed grant overhead $g=2$. Under limited disclosure and timely local recognition in two of four source types, $q=1/2$, so both owners strictly prefer preparation: A earns six against cost four, B earns eight against overhead two. Making the same usable grant known to the fee setter before the offer lowers A's gross return to three, despite all completed-setup types now being willing to trade. No preparation follows. If recognition is late in one of the two positive types, $q=1/4$ with no publication to B, so B switches to screening and A's gross return falls to $1.5$. A preapproved, binding fee of twelve restores preparation without choosing these disclosure powers if the constitution permits that commitment. Thus assignment flexibility is not necessary under a wider contract space.

\paragraph{Truthful evidence constrained by its physical suppliers.}
Now suppose some no-option cases can supply an authenticated failure receipt to the seller before its price is set, but the reporting installation must be chosen before the realized history. The option-holding cases never emit that failure receipt. Let absolute population mass $e$ be of reportable failures selected to send early; the residual class has mass $M=1-e$ and posterior $q/M$. Let $\eta$ be the additional fixed real cost to B of installing those early report paths. In the regime $qa<c$ and $qV<a$, neither no disclosure nor full truthful disclosure induces A to prepare. An installed set of sources restores mutually preferable preparation precisely when
\[
  c<aM\le qV,\qquad V-aM>g+\eta.
\]
The first two inequalities are A's investment and B's later pooling decisions; the last concerns B's ex ante participation in the install decision. Remaining types still satisfy the declared individual participation and can refuse. At the pooling threshold the displayed weak inequality assumes the announced price tie-break; use the strict inequality for robustness to either optimal offer.

For $q=1/4$, $c=5/2$, and other values as above, one reported failure group of mass $1/2$ leaves $M=1/2$. The silent category then pools; A receives three gross, and B eleven before grant and channel cost. Installing the early reporter for cost one is individually preferable to no preparation for both. Reporting only the other failure group of mass $1/4$ still screens; reporting all failures leaves A at $1.5$. If the two failure groups share an indivisible early reporting path, feasible values of $e$ are only $0$ and $3/4$, and neither works. Belief mixtures permitted by abstract persuasion can therefore be unattainable for actual authority and channel configurations. With finitely many disjoint indivisible evidence groups, the exact installed subset problem is a band constraint on achievable mass and a cost threshold per mass; the replay uses subset enumeration and a matched dynamic programme. Disjoint support and enforcement are necessary for that additive reduction; overlapping emitters and endogenously shifted failure probabilities change the finite model.

An early failure receipt volunteered by its beneficiary is not equivalent to a binding physical reporting rule. If A can choose after seeing its type to suppress a true negative receipt, any no-option type expected to disclose strictly prefers silence whenever silence has a positive probability of a pooling offer: disclosure leaves zero, silence offers a strictly positive rent. For $qV<a$, an agreement requiring those very types to disclose cannot be a sustainable equilibrium of this no-commitment game. If no such type discloses and high types cannot lawfully send a negative receipt, silence has the prior and B screens. A retained source publisher outside A's later discretion, an enforceable clause, or compensation for disclosure changes this result. The example stipulates an installed physical publisher with preassigned output and available report path; it does not derive that publisher from consent alone.

\paragraph{Joint law versus implementation.}
Consider another indistinguishable source pair where consent and price are valid in both histories, but in one an old use still physically occupies the joint apparatus. A controller seeing the same permitted signal cannot safely launch in the clear history: that would also launch in the occupied one. The robust rule is to start only in a cell where both recognized permission and actual clearance hold for every history. A distinct timely clearance witness permits the clear case, not the unsafe case. This pair is independent of the pricing model and does not replace its information conditions. Emergency stop and independent required safety work remain available; refusal is not counted as completing the optional shared service.

\paragraph{Comparators and verification.}
\href{https://www.cs.princeton.edu/courses/archive/spring09/cos444/papers/myerson81.pdf}{Myerson}'s incentive-compatible design and \href{https://doi.org/10.1016/j.jet.2016.10.008}{Hwang and Li}'s transparency of outside options are the closest pricing comparisons; the former already covers the basic optimal price argument and the latter studies bargaining with stochastically arriving options over multiple offers. \href{https://web.stanford.edu/~gentzkow/research/BayesianPersuasion.pdf}{Kamenica and Gentzkow}'s commitment to an information structure explains the unconstrained disclosure comparator; Bayesian plausibility alone does not supply a source or its publication power. \href{https://pages.nyu.edu/debraj/Papers/KonishiRay.pdf}{Konishi and Ray}'s dynamic coalition process is a broader state-and-continuation framework, not a substitute for this private-information fee game and timed effect guards. The strong matched direct alternatives used here are finite backward induction, generic two-type incentive-compatible optimization, and the reachable report-group optimization. They agree. No new general bargaining or access-control theorem, construction for faulty signers, independent transfer, or acceptance of R5 is claimed. There are no empirical estimates of payoff, fault masses or setup costs.

Replays: 9,216 physical comparisons, 4,624 exact price trees, 27,875 price checks after permitted reporting, and 512 separate local clearance comparisons. A generic incentive-compatible numerical linear programme agrees on 245 controls; its exact upper bound and selected inequalities check in Lean without sorryAx. The tests do not make an economic pricing envelope into a safety guarantee.
\end{document}
